\subsection*{Ordinary source mass and a full operator residual}
The Sobolev G1 estimate and the polynomial Fourier tail also give an
absolute ordinary-norm operator residual for this fixed source. This supplies the
Rayleigh and residual prerequisites in the spectral-overlap criterion;
it does not supply the overlap itself.

\begin{lemma}[A nonzero ordinary-norm limit of the repaired proxy]
\label{lem:v115-source-mass}
Put
\[
 h_\infty(u)=\frac{4\pi}{\sqrt3}u^2(2\pi u^2-3)e^{-\pi u^2},
 \qquad f_\infty=\mathcal E(h_\infty).
\]
Extend $p_\lambda$ by zero outside $I_\lambda$. Then
\begin{equation}\label{eq:v115-source-limit}
 \norm{p_\lambda-f_\infty}_{L^2(du/u)}=O(\lambda^{-1/2}),
 \qquad Jf_\infty=f_\infty.
\end{equation}
In particular $0<\norm{f_\infty}<\infty$. For
$N=N_\lambda=\lceil\lambda^8(1+L)\rceil$ and
$k_\lambda=P_Np_\lambda$, both $\norm{p_\lambda}$ and
$\norm{k_\lambda}$ tend to $\norm{f_\infty}$ and are bounded above
and away from zero. A valid eventual common lower bound is
\begin{equation}\label{eq:v115-mass-lower}
 m_*:=\frac14\left(\int_1^2h_\infty(u)^2\,du\right)^{1/2}>0.
\end{equation}
No numerical threshold in $\lambda$ is asserted.
\end{lemma}
\begin{proof}
The uniform fixed-mode Hermite approximation in
\cite[Lemma~7.2]{ConnesConsaniMoscovici2025} gives suitably normalized
$q_{j,\lambda}$ with
$\sup_{|u|\le\lambda}|q_{j,\lambda}(u)-H_j(u)|=O(\lambda^{-2})$,
$j=0,4$, where $H_j$ has unit $L^2(\mathbb R)$ norm. This rate survives
the manuscript's exact unit normalization: expanding the squared norm,
the cross term is $O(\lambda^{-2})\norm{H_j}_1$, the squared-error term
is $O(\lambda^{-3})$, and the omitted Hermite tail is Gaussian.
Thus $\norm{q_{j,\lambda}}=1+O(\lambda^{-2})$.

The limits of the coefficients in \eqref{eq:v12-coefficients} are
$-H_4(0)=-\sqrt3/(2\,2^{1/4})$ and $H_0(0)=2^{1/4}$, with error
$O(\lambda^{-2})$. Substitution of the normalized Hermite polynomials
gives exactly the displayed $h_\infty$, including its scalar factor.
The fixed repair has exponentially small amplitude by
\eqref{eq:v12-zero-value}. Consequently
\[
 \sup_{|u|\le\lambda}|\widetilde h_\lambda(u)-h_\infty(u)|
       \le C\lambda^{-2}.
\]
The limiting source is a Fourier-invariant combination of orders $0,4$,
with zero value at zero and zero integral. Poisson summation therefore
gives $Jf_\infty=f_\infty$. Its Gaussian decay at infinity, followed by
inversion at zero, gives square integrability.

On the physical window there are at most $\lambda/u$ compact-source
summands. Their accumulated approximation error is at most
$C\lambda^{-1}u^{-1/2}$. Its squared Haar norm is bounded by
\[
 C\lambda^{-2}\int_{1/\lambda}^\lambda u^{-2}\,du
       \le C\lambda^{-1}.
\]
The omitted Gaussian summands $nu>\lambda$ contribute an exponentially
smaller norm: bound their decreasing polynomial-Gaussian tail by its
first term plus $u^{-1}$ times its tail integral. The same estimate and
inversion control $f_\infty$ outside $I_\lambda$. Inversion averaging
is an $L^2$ contraction, proving \eqref{eq:v115-source-limit}.
For $u\ge1$, all summands $h_\infty(nu)$ are positive. Thus
$\norm{f_\infty}^2\ge\int_1^2h_\infty(u)^2\,du>0$.
Finally \eqref{eq:v114-high-tail} gives
\[
 \norm{(I-P_N)p_\lambda}
       \le C|B_\lambda|\lambda\sqrt{L/N}\longrightarrow0,
\]
so the projected vectors have the same nonzero norm limit. Choosing a
common sufficiently large threshold proves \eqref{eq:v115-mass-lower}.
\end{proof}

\begin{proposition}[A full residual after ordinary norm normalization]
\label{prop:v115-absolute-residual}
Use the same $N_\lambda,k_\lambda$ and the canonical closed operator
$\mathsf W_\lambda$ of Proposition~\ref{prop:v114-weil-core}.
For all sufficiently large $\lambda$, with $c=2\pi\lambda^2$,
\begin{equation}\label{eq:v115-absolute-residual}
 \norm{\mathsf W_\lambda p_\lambda}
 +\norm{\mathsf W_\lambda k_\lambda}
 +\norm{W_{\lambda,N_\lambda}k_\lambda}
       \le C\lambda^6 e^{-c/3}.
\end{equation}
For either source normalized to ordinary norm one, its Rayleigh value
and its full centered residual tend to zero at this rate. The finite
matrix assertion uses that same normalization and the original physical
Fourier cutoff.
\end{proposition}
\begin{proof}
All source vectors belong to $\mathcal D(\mathsf W_\lambda)$ by
Proposition~\ref{prop:v114-weil-core}. The weighted Chebyshev estimate
$P_\lambda=O(\lambda)$, together with $h_L=O(\lambda)$, gives
$\norm{[M_b,H]}=O(\lambda)$ in \eqref{eq:v114-b-bound}.
Corollary~\ref{cor:g1-complete-matrix}, including its separately proved
diagonal, gives uniformly
\[
 |d_n|\le C\{\lambda+1+\log(2+|n|/L)\}.
\]
Let $N=N_\lambda$ and $e=(I-P_N)p_\lambda$.
Using \eqref{eq:v114-high-tail} and comparing
$\sum_{n>N}\log^2(2+n/L)/n^2$ with its integral yields
\begin{align}
 \norm{\mathsf W_\lambda e}
 &\le C|B_\lambda|\lambda\sqrt{L/N}
        \{\lambda+1+\log(2+N/L)\}\notag\\
 &\le C|B_\lambda|\lambda^{-2}.
 \label{eq:v115-uniform-graph-tail}
\end{align}
Both the diagonal part and the bounded commutator are retained.

Next split the \emph{output} at an auxiliary integer $M>2N$.
This does not change the candidate or its finite matrix. For any
$v=P_Mv$ with $\norm v=1$,
\[
 \norm{D_{\log}v}\le2\pi M/L,\qquad
 |v(-a)|+|v(a)|\le2\sqrt{(2M+1)/L}.
\]
The full Sobolev estimate \eqref{eq:G1-sobolev}, whose constant is
independent of its test function, therefore implies
\begin{equation}\label{eq:v115-low-output}
 \norm{P_M\mathsf W_\lambda p_\lambda}
 \le C|B_\lambda|\mathfrak r(\lambda)
       \{1+M/L+\sqrt{M/L}\}.
\end{equation}
This includes the zero Fourier mode. It uses the already established
compact-source radicality and centered pole--prime cancellation, with
both source moments and inversion averaging intact. No fixed-band
constant is extrapolated to a growing band.

The diagonal has no cross block between $E_N$ and $E_M^\perp$.
For the complete off-diagonal matrix, Lemma~\ref{lem:g1-high-low}
gives the Hilbert--Schmidt bound
\[
 \norm{(I-P_M)\mathsf W_\lambda P_N}^2
 \le C\lambda^2\sum_{|m|\le N}\sum_{|n|>M}\frac1{(n-m)^2}
 \le C\lambda^2N/M.
\]
Lemma~\ref{lem:v115-source-mass} and
\eqref{eq:v115-uniform-graph-tail} now give
\begin{equation}\label{eq:v115-high-output}
 \norm{(I-P_M)\mathsf W_\lambda p_\lambda}
 \le C\lambda^5\sqrt{1+L}\,M^{-1/2}
         +C|B_\lambda|\lambda^{-2}.
\end{equation}
Thus all prime, pole and archimedean pieces remain in the actual matrix;
no divergent exterior pole or prime norm sum is estimated separately.

Take $M=\lceil e^{2c/3}\rceil$, which exceeds $2N$ eventually.
By \eqref{eq:v12-combined-endpoint} and
\eqref{eq:v13-endpoint-comparison},
$|B_\lambda|\le C\lambda^{11/2}e^{-c}$, while
$\mathfrak r(\lambda)$ is bounded. Equations
\eqref{eq:v115-low-output}--\eqref{eq:v115-high-output} are therefore
at most $C\lambda^6e^{-c/3}$ in total. This proves the first term of
\eqref{eq:v115-absolute-residual};
\eqref{eq:v115-uniform-graph-tail} proves the second, and compression
contractivity proves the third.
Finally ordinary normalization is legitimate by
\eqref{eq:v115-mass-lower}. For a unit vector $u$, letting
$\alpha=\ip{\mathsf W_\lambda u}u$, one has
$|\alpha|\le\norm{\mathsf W_\lambda u}$ and
$\norm{(\mathsf W_\lambda-\alpha)u}\le2\norm{\mathsf W_\lambda u}$.
The same argument applies to the compressed matrix.
\end{proof}

The preceding proposition supplies no endpoint-normalized operator residual:
dividing its bound by $|B_\lambda|$ leaves an exponentially growing
upper bound. Nor does it identify the spectral point near zero as the
lowest one. The elementary family $\operatorname{diag}(-1,\epsilon)$,
with candidate $(0,1)$, has Rayleigh value $\epsilon$ and centered
residual zero while its lowest level stays at $-1$. An independent
quantitative ground-space overlap or signed complementary lower bound
is still necessary. No estimate on the corrected sampler's graph defect,
or on a growing family of prolate modes, is inferred here.
